\chapter{Experimental Design}
\label{ch:design}

\section{Data, training, and the matched-budget protocol}

Both arms train on the PDBbind general set \citep{liu2017pdbbind} and are
evaluated on the PoseBusters benchmark \citep{buttenschoen2024posebusters},
$209$ complexes after the filtering inherited from the parent codebase. Splits,
fragmentation, conformer generation and featurisation are shared, so the two
arms consume byte-identical inputs.

\emph{Matched budget} here means identical wall-clock allocation on identical
hardware, identical batch size, identical optimiser and identical annealing
schedule --- not identical loss values or identical numbers of epochs, which
are outcomes rather than controls. Table~\ref{tab:configs} records the
configurations.

\begin{table}[htbp]
\centering
\caption{Training configurations. The $12$\,h pair is the matched comparison
available at the time of writing; the long pair is the primary experiment.}
\label{tab:configs}
\small
\begin{tabular}{@{}lccccc@{}}
\toprule
Run & Arm & Wall-clock & Hardware & Batch & Status \\
\midrule
Early pair & SigmaFlow & $12$\,h & $1\times$L40S & 8 & complete \\
Early pair & SigmaDock  & $12$\,h & $1\times$L40S & 8 & complete \\
Long pair  & SigmaFlow & $72$\,h & $1\times$L40S & \tbd & \dnote{RESULT PENDING}{queued} \\
Long pair  & SigmaDock  & $72$\,h & $1\times$L40S & \tbd & \dnote{RESULT PENDING}{queued} \\
Source screen & SigmaFlow-informative & $6$\,h & $1\times$L40S & \tbd & \dnote{RESULT PENDING}{gated} \\
Source screen & SigmaFlow-Haar control & $6$\,h & $1\times$L40S & \tbd & \dnote{RESULT PENDING}{gated} \\
\bottomrule
\end{tabular}
\end{table}

For calibration against the literature: SigmaDock's published results were
obtained with roughly $384$ GPU-hours, whereas the early pair here used about
$12$ each. This thesis therefore operates at a small fraction of the compute
behind the published numbers, and absolute accuracies are correspondingly low
for \emph{both} arms. The comparison remains valid because the budget is
matched; the benchmark comparison is not attempted.

\section{Trajectory snapshots are not independently annealed endpoints}
\label{sec:snapshot-caveat}

The long runs checkpoint at fixed wall-clock intervals, and it is tempting to
read the resulting sequence as performance ``at $12$\,h, at $24$\,h, and so
on''. That reading is wrong and the distinction is maintained throughout
Chapter~\ref{ch:results}.

The learning rate follows a cosine schedule whose length is set by the total
step budget. A snapshot taken part-way through a $72$\,h schedule is therefore
a model \emph{mid-anneal}, at a learning rate far above its terminal value,
and it is not the model one would obtain by training for that shorter budget
with a schedule annealed to it. Two distinct objects follow:

\begin{itemize}[leftmargin=*,itemsep=2pt]
\item a \textbf{training trajectory} --- metrics against optimisation step for
      a \emph{fixed} schedule, with both arms sharing that schedule. This is a
      controlled comparison and it answers whether the rotational deficit
      narrows as training proceeds;
\item \textbf{annealed endpoints} --- separately scheduled runs, each having
      completed its own anneal. The $12$\,h pair and the $72$\,h endpoint are
      of this kind.
\end{itemize}

Trajectories are plotted as curves; annealed endpoints are plotted as isolated
markers and never joined to them. No claim of the form ``performance scales
with compute as \ldots'' is made, since establishing that would require a
separate fully annealed run per budget.

\section{Metrics, nulls, and statistical treatment}
\label{sec:stats}

Primary quantities are chosen for low sampling variance, so that a trajectory
is readable rather than noisy:

\begin{itemize}[leftmargin=*,itemsep=1pt]
\item median per-fragment rotation error, against the Haar reference
      \eqref{eq:haar-null}, and the enrichment below $\Ang{30}$;
\item median fragment centroid displacement;
\item median RMSD, raw and after optimal superposition;
\item validation loss, separately for the rotational and translational terms;
\item the diagnostic ratios of \S\ref{sec:predictions}: predicted-to-target
      magnitude ratio, and predicted-target cosine stratified by $t$.
\end{itemize}

Threshold success rates --- the fraction below $2$\,\AA{} --- are reported in
tables with confidence intervals but are \emph{not} plotted over training. At
the rates observed so far and $n = 209$, their standard error is of the same
order as any plausible change, so a curve would display sampling noise.

All paired comparisons are computed per complex, with bootstrap confidence
intervals over complexes. Where several step counts or variants are compared
against a common reference, the multiplicity is stated and no individual
comparison is claimed significant on the strength of an uncorrected interval.
Stochasticity across sampling seeds is reported separately from stochasticity
across complexes, since the two have been observed to be of comparable size
and conflating them would overstate precision.
